%! Logarithmic Functions — Unit 5, Lesson 5.3 — Guided Notes
\documentclass[10pt]{article}
\usepackage{precalculus-article}
\usepackage{precalculus-boxes}
\newcommand{\TallMath}[1]{$\displaystyle #1\rule[-1.4em]{0pt}{3.2em}$}

\begin{document}

\pageheader{Unit 5, Lesson 5.3}{Guided Notes}
\namedateperiod

\vspace{0.08in}

\begin{objectivebox}
\textbf{LO 2.11.A:} Identify key characteristics of logarithmic functions.

By the end of this lesson I will be able to:
\begin{itemize}[leftmargin=1.5em,itemsep=3pt,topsep=3pt]
  \item State the domain, range, and vertical asymptote of $f(x) = a\log_b x$.
  \item Describe the end behavior, monotonicity, and concavity of a logarithmic function.
  \item Use the additive transformation identification test to decide if a function is logarithmic.
\end{itemize}
\end{objectivebox}

\vspace{0.06in}

\begin{vocabbox}
\termblanklong{Logarithmic function (general form)}
\termblanklong{Domain of a logarithmic function}
\termblanklong{Vertical asymptote}
\termblanklong{End behavior}
\termblanklong{Monotone increasing / decreasing}
\termblanklong{Concave up / concave down}
\termblanklong{Additive transformation identification test}
\end{vocabbox}

\vspace{0.06in}

\begin{hookbox}
\textbf{Richter Scale Context:} The Richter scale is defined by $M = \log_{10}(I)$, where
$I$ is earthquake intensity.

\vspace{4pt}
If intensity $I = 10$, what is the magnitude? \blank{2cm}
If intensity $I = 100$, magnitude $= $ \blank{2cm}
If intensity $I = 1000$, magnitude $= $ \blank{2cm}

\vspace{4pt}
Each time $I$ multiplies by \blank{1.5cm}, the magnitude $M$ increases by \blank{1.5cm}.

\vspace{4pt}
What inputs $I$ are physically meaningful for this function? \writeline
\end{hookbox}

\vspace{0.06in}

\begin{notesbox}{Domain and Range (EK 2.11.A.1)}
For $f(x) = a\log_b x$ in general form:

\vspace{6pt}
\begin{center}
\renewcommand{\arraystretch}{1.8}
\begin{tabular}{|l|l|}
\hline
\rowcolor{sky}
\textbf{Characteristic} & \textbf{Value} \\
\hline
Domain & \blank{5cm} \\
\hline
Range  & \blank{5cm} \\
\hline
Vertical asymptote & $x = $ \blank{2cm} \\
\hline
\end{tabular}
\end{center}

\vspace{6pt}
Why can't we evaluate $\log_b 0$ or $\log_b(-3)$? \writeline
\end{notesbox}

\vspace{0.06in}

\begin{notesbox}{Characteristics from the Inverse (EK 2.11.A.2)}
Because $f(x) = \log_b x$ is the inverse of $g(x) = b^x$, it inherits structural properties.
Complete the comparison table.

\vspace{6pt}
\begin{center}
\renewcommand{\arraystretch}{1.8}
\begin{tabular}{|l|l|l|}
\hline
\rowcolor{sky}
\textbf{Property} & \textbf{Exponential $g(x) = b^x$} & \textbf{Logarithmic $f(x) = \log_b x$} \\
\hline
Always increasing ($b>1$)? & Yes & \blank{3cm} \\
\hline
Always decreasing ($0<b<1$)? & Yes & \blank{3cm} \\
\hline
Has local extrema (natural domain)? & No & \blank{3cm} \\
\hline
Has an inflection point? & No & \blank{3cm} \\
\hline
\end{tabular}
\end{center}

\vspace{6pt}
So logarithmic functions are always \blank{3cm} or always \blank{3cm}, and their graphs
are either always \blank{3cm} or always \blank{3cm}.
\end{notesbox}

\vspace{0.06in}

\begin{notesbox}{End Behavior (EK 2.11.A.4)}
Fill in the end behavior for each case. The sign of $a$ and the value of $b$ both matter.

\vspace{6pt}
\begin{center}
\renewcommand{\arraystretch}{2.0}
\begin{tabular}{|l|l|l|}
\hline
\rowcolor{sky}
\textbf{Conditions} & \TallMath{\lim_{x \to 0^+} a\log_b x} & \TallMath{\lim_{x \to \infty} a\log_b x} \\
\hline
$b > 1$, $a > 0$ & \blank{2.5cm} & \blank{2.5cm} \\
\hline
$b > 1$, $a < 0$ & \blank{2.5cm} & \blank{2.5cm} \\
\hline
$0 < b < 1$, $a > 0$ & \blank{2.5cm} & \blank{2.5cm} \\
\hline
\end{tabular}
\end{center}
\end{notesbox}

\vspace{0.06in}

\begin{notesbox}{Additive Transformation Identification Test (EK 2.11.A.3)}
\textbf{Key Idea:} If $g(x) = f(x + k)$ for $k \ne 0$, then $g$ is an additive
transformation of $f$.

\begin{itemize}[leftmargin=1.5em, itemsep=4pt, topsep=4pt]
  \item If the \emph{inputs} of $g(x) = f(x+k)$ are proportional over equal-length output
        intervals, then $f$ is logarithmic.
  \item If the inputs of $f$ itself are proportional over equal-length output intervals,
        that confirms $f$ is logarithmic.
\end{itemize}

\vspace{6pt}
\textit{Example.} Below is a table for $f(x) = \log_3 x$. Check: as the output increases
by 1, does the input multiply by a constant?

\vspace{4pt}
\begin{center}
\renewcommand{\arraystretch}{1.8}
\begin{tabular}{|c|c|c|}
\hline
\rowcolor{sky}
$x$ (input) & $f(x)$ (output) & Ratio of consecutive inputs \\
\hline
$1$  & $0$ & --- \\
\hline
$3$  & $1$ & $3/1 = $ \blank{1.5cm} \\
\hline
$9$  & $2$ & $9/3 = $ \blank{1.5cm} \\
\hline
$27$ & $3$ & $27/9 = $ \blank{1.5cm} \\
\hline
$81$ & $4$ & $81/27 = $ \blank{1.5cm} \\
\hline
\end{tabular}
\end{center}

\vspace{6pt}
The inputs are proportional (each multiplied by \blank{1.5cm}), so $f$ \blank{3cm}
(is / is not) logarithmic.
\end{notesbox}

\end{document}
